\documentclass{article}
\usepackage[utf8]{inputenc}
\usepackage[T1]{fontenc}
\usepackage{geometry}
\usepackage{graphicx}
\usepackage{float}
\usepackage{hyperref}
\usepackage[english,indonesian]{babel}

\title{Optimasi Pemrosesan Template LaTeX Menggunakan Pendekatan Berbasis AI}
\author{Academic Support Agent}
\date{\today}

\begin{document}
\maketitle

\begin{abstract}

\end{abstract}

\chapter{PENDAHULUAN}


\chapter{TINJAUAN PUSTAKA}


% Original content (commented out):
% % Inicio del preámbiulo

\documentclass[letterpaper,12pt]{article} %Modifica el tipo de documento y el tamaño de la letra.
\usepackage[utf8]{inputenc} %Formato UTF-8 para caracteres especiales.
\usepackage[shortlabels]{enumitem}
\usepackage[spanish,mexico]{babel} 
\usepackage{amsmath,amssymb,amsfonts,latexsym,cancel}
\usepackage{hyperref}
\usepackage{wrapfig}
\usepackage[rflt]{floatflt}
\usepackage[pdftex]{graphicx}
\usepackage{fancyhdr} %Paquete para el header y el formato de la portada. No sugiero borrarlo!
\usepackage{float}
\usepackage{longtable,multirow,booktabs}
\usepackage{cite}
\usepackage{wrapfig}
% \usepackage[square,numbers]{natbib}
\usepackage{multicol}
\usepackage{caption}
\usepackage[]{sidecap}
\usepackage{adjustbox}
\usepackage{parskip}
\usepackage{enumitem}
\usepackage{tikz}
\usepackage{lipsum}
\usepackage[]{xcolor}

%Fin de Préambulo

% Variables del documento
% Document Variables 
\newcommand{\myMateria}{Curso}
\newcommand{\myGrupo}{1}
\newcommand{\mySemester}{2022-1}
\newcommand{\MyReport}{Titulo del reporte}
\newcommand{\myUnidad}{1}
\newcommand{\myDate}{8 de abril de 2022}
\newcommand{\myName}{Darío Ramsés Gutiérrez Rodríguez}
\newcommand{\myTeacher}{Isaac Ramírez}



%Inicio formato de Página. 

\textheight = 21cm %Medidas de la  página
\textwidth = 18cm  %Medidas de la página
\topmargin = -2cm  %Medidas de la página    
\oddsidemargin = -0.8cm %Medidas de la página
\pagestyle{fancy} %Diseño de la página

\fancyhf{}
\lhead{\myMateria}%%LeftHead
\chead{\includegraphics[ height=1cm]{Imágenes/logo-itcr.png}}%%CenterHead
%\lfoot{USM}
% \rhead{Unidad \myUnidad}%%RightHead

\setlength{\columnsep}{4mm}%Comandos para el formato de la página.
%\setlength{\parindent}{4em}%Sangría al comenzar un nuevo párrafo.
\setlength{\parindent}{0.5in}
%\setlength{\parindent}{4em}%Sangría al comenzar un nuevo párrafo.
\setlength{\parskip}{1em}%Distancia entre párrafos.
\renewcommand{\baselinestretch}{1.0}% Espacio entre línea y línea o interlineado.
\setlength{\headheight}{33pt}
\fancyfoot[C,CO]{\thepage} %Logo de LaTeX y pie de página.

%Fin formato de Página

%Aquí inicia el documento.

\begin{document}

    %LaTeX te hace el índice automáticamente conforme añades secciones en tu documento.
    \input{portada.tex}
    \newpage
    \tableofcontents
    \newpage
%Inicio parte opcional. Esta parte la puedes quitar si deseas, es por si te piden formatos para
%evidencias de certificación de los laboratorios con números de cuenta o te piden abstracts en tus %trabajos.

\title{\myMateria \\\textbf{\MyReport} \\ } 

\author{ \normalsize{\texttt{\myName}} }
\date{\myDate}
\maketitle
\thispagestyle{fancy}

\begin{abstract}
En esta seccion va el resumen del proyecto
\end{abstract}

%Fin parte opcional

\input{contenido}

\nocite{*}
\bibliographystyle{ieeetr}
\bibliography{Libros.bib}


%Pequeño disclaimer por si usas fotos con derechos de autor que sacaste de Internet. Nunca está de más. 
\centering \small Los créditos de las fotografías pertenecen a sus respectivos autores.

\centering\vspace*{\fill} 
\end{document}

 %Fin del documento.

\end{document}